\section{Transformaci\'on digital y agricultura inteligente}

\subsection{Agricultura 4.0}

La Agricultura 4.0 integra sensores, conectividad, automatizaci\'on y anal\'itica de datos para optimizar decisiones en tiempo real \citep{Chen2023}. En cultivos de alto valor como la palta Hass, la digitalizaci\'on abarca desde monitoreo de suelo y clima hasta trazabilidad postcosecha y anal\'itica comercial. Aunque La Libertad cuenta con fundos tecnificados, persiste una brecha entre grandes exportadores y peque\~nos productores en adopci\'on de herramientas digitales.

La transformaci\'on digital no se limita a la productividad agron\'omica: incluye la capacidad de procesar informaci\'on de mercado, simular escenarios econ\'omicos y comunicar recomendaciones oportunas. \citet{Pantaleon2024} se\~nala que la sostenibilidad de exportadores peruanos depende cada vez m\'as de gesti\'on de riesgo, cumplimiento normativo y respuesta r\'apida a se\~nales de demanda, \'ambitos donde la anal\'itica de datos aporta ventajas competitivas.

Para productores de palta, Agriculture 4.0 implica combinar datos locales (rendimiento, costos, calendario de cosecha) con datos globales (precios FOB, vol\'umenes de exportaci\'on, tipo de cambio). Sin esta integraci\'on, la tecnificaci\'on productiva no se traduce necesariamente en mejores decisiones de inversi\'on a mediano plazo.

\subsection{Uso de Big Data e Inteligencia Artificial}

El volumen creciente de datos comerciales, aduaneros y agroclim\'aticos habilita enfoques de Big Data e Inteligencia Artificial (IA) aplicados a la agricultura \citep{Chen2023, Murugesan2022, ReisFilho2022}. Fuentes como Aduanet, INEI, FAOSTAT y bases sectoriales permiten construir series temporales de precio, volumen y destino con granularidad mensual o semanal \citep{UNDataFAO}.

\citet{Chen2023} demuestran modelos h\'ibridos que combinan descomposici\'on de se\~nales (VMD) con redes LSTM para pronosticar precios agr\'icolas, superando enfoques univariados tradicionales. \citet{Murugesan2022} comparan arquitecturas LSTM, BiLSTM, CNN-LSTM y ConvLSTM en commodities agr\'icolas, evidenciando ganancias de precisi\'on cuando se modelan dependencias temporales complejas. \citet{ReisFilho2022} enriquecen series con datos textuales ---noticias, reportes--- para capturar choques no reflejados inmediatamente en precios hist\'oricos.

En el contexto peruano, la IA puede procesar patrones estacionales de exportaci\'on peruana, correlaciones con or\'igenes competidores y rezagos de precio para generar pron\'osticos FOB. No obstante, la calidad del dato ---consistencia aduanera, tratamiento de outliers, alineaci\'on temporal--- condiciona la utilidad de cualquier modelo avanzado.

\subsection{Beneficios de la anal\'itica predictiva para la agricultura}

La anal\'itica predictiva transforma datos hist\'oricos y explicativos en proyecciones accionables: precios esperados, rangos de incertidumbre, escenarios de rentabilidad y alertas tempranas \citep{Theofilou2025, Xia2024}. Para productores de palta, los beneficios incluyen mejor timing de decisiones de siembra, reducci\'on de exposici\'on a a\~nos de precios deprimidos y mayor transparencia frente a intermediarios.

\citet{Xia2024} comparan ARIMA y LSTM en productos agr\'icolas, concluyendo que modelos de aprendizaje profundo capturan mejor no linealidades cuando hay suficiente historial. \citet{Velasquez2025} vinculan cuantificaci\'on de riesgo con herramientas de apoyo a exportadores, enfoque extensible a productores mediante simuladores simplificados. En la presente investigaci\'on se plantean cinco soluciones anal\'iticas ---rentabilidad, escenarios, predicci\'on FOB, hist\'orico interactivo y reportes de ciclos--- que articulan predicciones con decisiones de negocio.

La clave no es solo predecir con bajo error (RMSE, MAPE), sino comunicar resultados en lenguaje comprensible: ``si el precio FOB cae 10\%, el margen proyectado ser\'ia X''. Este puente entre modelo y usuario es central en agricultura inteligente orientada a peque\~nos productores, no solo a analistas de trading agr\'icola.
